\sectionkicker{Source-backed technical notes}
\subsection*{AgentはUIを出て実行系になる — Harness・Context・Tool・Browser/Computer Use: Technical Notes}
本欄は記事本文で圧縮した一次資料上の情報を、比較・再検証しやすい形で整理したものである。時系列、確認済みの主張、留意点、一次資料URLを掲載する。
\medskip
\subsection*{Theme at a glance}
\addcontentsline{toc}{subsection}{Theme at a glance}
\begin{center}
\footnotesize
\begin{tabularx}{\linewidth}{@{}p{0.20\linewidth}p{0.10\linewidth}p{0.14\linewidth}X@{}}
\toprule
資料 & 位置づけ & 種別 & 時系列 \\
\midrule
ChatGPT agent & 主要資料 & Agent & 2025-07-17 (Agent（公開）); 2025-07-17 (System Card公開)      \\
Claude 4.5 model and agent platform sequence & 主要資料 & モデル更新 & 2025-09-29 (Claude Sonnet 4.5（公開）); 2025-10-15 (Claude Haiku 4.5（公開）); 2025-11-24 (Claude Opus 4.5（公開）)      \\
Gemini API H2 2025 lifecycle & 主要資料 & API & 2025-07-07 (Batch Mode（公開）); 2025-07-17 (Veo 3（Preview）); 2025-07-22 (Gemini 2.5 Flash-Lite（公開）); 2025-08-14 (Imagen 4（一般提供）); 2025-08-26 (Gemini 2.5 Image（Preview）); 2025-09-09 (Veo 3（一般提供）); 2025-09-25 (Robotics-ER 1.5（Preview）); 2025-10-07 (Computer Use（Preview）); 2025-10-15 (Veo 3.1（Preview）); 2025-11-18 (Gemini 3 Pro（Preview）); 2025-12-11 (Interactions API / Deep Research（Preview）); 2025-12-17 (Gemini 3 Flash（Preview）)      \\
AgentKit, Evals, and RFT for agents & 補足資料 & Framework & 2025-10-06 (製品ツール（公開）)      \\
Codex September upgrades & 補足資料 & 製品 & 2025-09-15 (Coding Agent（更新）)      \\
\bottomrule
\end{tabularx}
\end{center}
\normalsize

\begin{technicalnote}{ChatGPT agent}{主要資料}
\begingroup
% reader-facing Technical Notes break policy
\widowpenalty=10000
\clubpenalty=10000
\displaywidowpenalty=10000
\begin{tabularx}{\linewidth}{@{}>{\bfseries}p{0.22\linewidth}X@{}}
組織 & OpenAI \\
種別 & Agent \\
時系列 & 2025-07-17 (Agent（公開）); 2025-07-17 (System Card公開) \\
\end{tabularx}
\smallskip
{\bfseries 一次資料から整理したtechnical points}
\begin{itemize}[leftmargin=1.5em,itemsep=0.35em]
\item \textbf{一次情報で確認できる事実}: OpenAIは7月17日、調査、ブラウザ操作、コード／ツール実行をユーザーの指示下で行うChatGPT agentを公開した。
\item \textbf{一次情報で確認できる事実}: 同日公開のSystem Cardは、agentic systemに対して評価したリスクとcontrolの境界を記録した。
\end{itemize}
\begin{minipage}{\linewidth}
% reader-facing Technical Notes generic boundary/source tail group
{\bfseries 読む際の境界}
\begin{itemize}[leftmargin=1.5em,itemsep=0.35em]
\item \textbf{分析上の留意点}: 「ChatGPT agent」の扱いでは、一次資料で確認できる範囲、提供時点、評価条件を維持し、記載範囲を超えて一般化しない。
\end{itemize}
\begin{samepage}
{\bfseries 一次資料}
\begingroup\sloppy
\Urlmuskip=0mu plus 2mu\relax
\begin{itemize}[leftmargin=1.5em,itemsep=0.25em]
\item {\scriptsize\url{https://openai.com/index/chatgpt-agent-system-card}}
\item {\scriptsize\url{https://openai.com/index/introducing-chatgpt-agent}}
\end{itemize}
\endgroup
\end{samepage}
\end{minipage}
\endgroup
\end{technicalnote}
\medskip

\begin{technicalnote}{Claude 4.5 model and agent platform sequence}{主要資料}
\begingroup
% reader-facing Technical Notes break policy
\widowpenalty=10000
\clubpenalty=10000
\displaywidowpenalty=10000
\begin{tabularx}{\linewidth}{@{}>{\bfseries}p{0.22\linewidth}X@{}}
組織 & Anthropic \\
種別 & モデル更新 \\
時系列 & 2025-09-29 (Claude Sonnet 4.5（公開）); 2025-10-15 (Claude Haiku 4.5（公開）); 2025-11-24 (Claude Opus 4.5（公開）) \\
\end{tabularx}
\smallskip
{\bfseries 一次資料から整理したtechnical points}
\begin{itemize}[leftmargin=1.5em,itemsep=0.35em]
\item \textbf{一次情報で確認できる事実}: Anthropicは9月29日、Claude Sonnet 4.5を、Claude Codeの更新およびClaude Agent SDKとあわせて公開した。
\item \textbf{一次情報で確認できる事実}: Anthropicは10月15日にClaude Haiku 4.5、11月24日にClaude Opus 4.5を公開し、いずれもAPIで提供した。
\item \textbf{分析上の留意点}: Anthropicの2025年後半の系列では、モデル公開がコーディング、Computer Use、コンテキスト・メモリ、エージェント構築の各利用面と結び付いて展開された。
\end{itemize}
\begin{minipage}{\linewidth}
% reader-facing Technical Notes generic boundary/source tail group
{\bfseries 読む際の境界}
\begin{itemize}[leftmargin=1.5em,itemsep=0.35em]
\item \textbf{分析上の留意点}: 保存したnewsroomのスナップショットはページ分割された継続更新型indexで、先頭HTMLは主に2026年の項目を示していた。2025年後半の年表は、個別のAnthropic公式ページを追跡して再構成した。
\end{itemize}
\begin{samepage}
{\bfseries 一次資料}
\begingroup\sloppy
\Urlmuskip=0mu plus 2mu\relax
\begin{itemize}[leftmargin=1.5em,itemsep=0.25em]
\item {\scriptsize\url{https://www.anthropic.com/news}}
\end{itemize}
\endgroup
\end{samepage}
\end{minipage}
\endgroup
\end{technicalnote}
\medskip

\begin{technicalnote}{Gemini API H2 2025 lifecycle}{主要資料}
\begingroup
% reader-facing Technical Notes break policy
\widowpenalty=10000
\clubpenalty=10000
\displaywidowpenalty=10000
\begin{tabularx}{\linewidth}{@{}>{\bfseries}p{0.22\linewidth}X@{}}
組織 & Google \\
種別 & API \\
時系列 & 2025-07-07 (Batch Mode（公開）); 2025-07-17 (Veo 3（Preview）); 2025-07-22 (Gemini 2.5 Flash-Lite（公開）); 2025-08-14 (Imagen 4（一般提供）); 2025-08-26 (Gemini 2.5 Image（Preview）); 2025-09-09 (Veo 3（一般提供）); 2025-09-25 (Robotics-ER 1.5（Preview）); 2025-10-07 (Computer Use（Preview）); 2025-10-15 (Veo 3.1（Preview）); 2025-11-18 (Gemini 3 Pro（Preview）); 2025-12-11 (Interactions API / Deep Research（Preview）); 2025-12-17 (Gemini 3 Flash（Preview）) \\
\end{tabularx}
\smallskip
{\bfseries 一次資料から整理したtechnical points}
\begin{itemize}[leftmargin=1.5em,itemsep=0.35em]
\item \textbf{一次情報で確認できる事実}: Gemini APIの変更履歴には、非同期バッチ処理、Embedding、動画、画像、ネイティブ音声、ロボティクス、Computer Use、Gemini 3、統合Interactions API、Deep Research agentにまたがる高密度な2025年後半の更新が記録されている。
\item \textbf{一次情報で確認できる事実}: 変更履歴はPreview、一般提供、非推奨化、停止を明示的に区別しており、これらの提供状態自体が技術年表の一部である。
\item \textbf{一次情報で確認できる事実}: Gemini 3 Pro Previewは11月18日にAPIへ追加され、Gemini 3 Flash Previewは12月17日に続いた。
\end{itemize}
\begin{minipage}{\linewidth}
% reader-facing Technical Notes generic boundary/source tail group
{\bfseries 読む際の境界}
\begin{itemize}[leftmargin=1.5em,itemsep=0.35em]
\item \textbf{分析上の留意点}: この変更履歴はサービス／APIの年表であり、上流の研究・モデル公開年表ではない。Previewから一般提供への移行や停止を、最初のモデル発表と同一イベントへ統合しない。
\end{itemize}
\begin{samepage}
{\bfseries 一次資料}
\begingroup\sloppy
\Urlmuskip=0mu plus 2mu\relax
\begin{itemize}[leftmargin=1.5em,itemsep=0.25em]
\item {\scriptsize\url{https://ai.google.dev/gemini-api/docs/changelog}}
\end{itemize}
\endgroup
\end{samepage}
\end{minipage}
\endgroup
\end{technicalnote}
\medskip

\begin{technicalnote}{AgentKit, Evals, and RFT for agents}{補足資料}
\begingroup
% reader-facing Technical Notes break policy
\widowpenalty=10000
\clubpenalty=10000
\displaywidowpenalty=10000
\begin{tabularx}{\linewidth}{@{}>{\bfseries}p{0.22\linewidth}X@{}}
組織 & OpenAI \\
種別 & Framework \\
時系列 & 2025-10-06 (製品ツール（公開）) \\
\end{tabularx}
\smallskip
{\bfseries 一次資料から整理したtechnical points}
\begin{itemize}[leftmargin=1.5em,itemsep=0.35em]
\item \textbf{一次情報で確認できる事実}: OpenAIはAgentKitを発表し、あわせてエージェント向け評価機能の拡充と強化学習によるファインチューニングを提示した。
\end{itemize}
\begin{minipage}{\linewidth}
% reader-facing Technical Notes generic boundary/source tail group
{\bfseries 読む際の境界}
\begin{itemize}[leftmargin=1.5em,itemsep=0.35em]
\item \textbf{分析上の留意点}: ここで確認できるのは一次資料に記録された範囲の事実であり、ベンダー・プロジェクト・著者による評価を独立再現された結果として扱うものではない。
\end{itemize}
\begin{samepage}
{\bfseries 一次資料}
\begingroup\sloppy
\Urlmuskip=0mu plus 2mu\relax
\begin{itemize}[leftmargin=1.5em,itemsep=0.25em]
\item {\scriptsize\url{https://openai.com/index/introducing-agentkit}}
\end{itemize}
\endgroup
\end{samepage}
\end{minipage}
\endgroup
\end{technicalnote}
\medskip

\begin{technicalnote}{Codex September upgrades}{補足資料}
\begingroup
% reader-facing Technical Notes break policy
\widowpenalty=10000
\clubpenalty=10000
\displaywidowpenalty=10000
\begin{tabularx}{\linewidth}{@{}>{\bfseries}p{0.22\linewidth}X@{}}
組織 & OpenAI \\
種別 & 製品 \\
時系列 & 2025-09-15 (Coding Agent（更新）) \\
\end{tabularx}
\smallskip
{\bfseries 一次資料から整理したtechnical points}
\begin{itemize}[leftmargin=1.5em,itemsep=0.35em]
\item \textbf{一次情報で確認できる事実}: OpenAIはCodexの更新を発表し、ターミナル、IDE、Web、モバイルへ利用面を広げるとともに、より独立したタスク実行を強調した。
\end{itemize}
\begin{minipage}{\linewidth}
% reader-facing Technical Notes generic boundary/source tail group
{\bfseries 読む際の境界}
\begin{itemize}[leftmargin=1.5em,itemsep=0.35em]
\item \textbf{分析上の留意点}: ここで確認できるのは一次資料に記録された範囲の事実であり、ベンダー・プロジェクト・著者による評価を独立再現された結果として扱うものではない。
\end{itemize}
\begin{samepage}
{\bfseries 一次資料}
\begingroup\sloppy
\Urlmuskip=0mu plus 2mu\relax
\begin{itemize}[leftmargin=1.5em,itemsep=0.25em]
\item {\scriptsize\url{https://openai.com/index/introducing-upgrades-to-codex}}
\end{itemize}
\endgroup
\end{samepage}
\end{minipage}
\endgroup
\end{technicalnote}
\medskip
